\Cref{thm:silence} says liars cannot remove information from a receiver that knows the joint law. It does not say how much information they add, or how much of it a label-free receiver can recover. We measure both with two oracles that see the gold labels of the HISTORY split. The \emph{liar-removal oracle} knows who lies, removes every liar, and decodes the honest agents with their true accuracies and true source groups. The \emph{known-channel oracle} keeps every agent, with true accuracies and true source groups. We call the gap between them the \emph{information budget} of an attack. E9 evaluates both oracles, \race\ and the extension below on \nWorldsExt\ worlds: 14 symbolic and 4 LLM attack settings, $f\in\{0.3,0.5,0.7\}$, six benchmarks and three seeds. \Cref{fig:budget} shows every method relative to the liar-removal oracle.

\paragraph{Truth-dependent liars are worth more than their absence.} For independent and gate-aware liars the budget is positive and grows with $f$. For independent liars it is \budgetIndepThree, \budgetIndepFive\ and \budgetIndepSeven\ points at $f=0.3$, $0.5$ and $0.7$. At $f=0.7$, three honest agents plus seven liars decode better than the three honest agents alone. \race\ captures most of this budget without labels and without knowing who lies. At $f=0.7$ it scores \extRaceIndepSeven\% against independent liars, while the oracle that removes every liar scores \extRemovalIndepSeven\%. Pooling tasks over benchmarks, \race\ is significantly \emph{above} the liar-removal oracle in \pooledRemovalWins\ of \pooledRemovalCells\ (attack, $f$) cells. All of these cells are truth-dependent attacks: independent, gate-aware with $p<1$ at every $f$, and the coherent bloc and attractor at $f=0.7$. A rule that only filters suspected liars sees at most the honest reports, so in expectation it cannot beat the Bayes receiver on those reports, for which this oracle is a proxy. Exceeding it requires reading the liars.

\paragraph{LLM deceivers: the value is neutralisation, not extra information.} For the four replayed LLM attacks the budget is close to zero: \budgetLlmThree, \budgetLlmFive\ and \budgetLlmSeven\ points on average. Deception prompts produce lies that are only weakly tied to the truth. At the one-coin level they carry little beyond what the honest agents already say. \race's contribution is then to cancel the damage. At $f=0.7$, majority vote falls to \extLlmMajSeven\% while \race\ scores \extLlmRaceSeven\%. That recovers \recoveryLlmSeven\% of the gap between the receiver alone (\extLlmSelfSeven\%) and the liar-removal oracle (\removalSeven\%); at $f=0.5$ it recovers \recoveryLlmFive\%. In the per-benchmark cells, \race\ is statistically indistinguishable from the liar-removal oracle in \removalLlmTies\ of \removalLlmCells\ LLM cells and below it in \removalLlmLosses.

\paragraph{A negative budget certifies a model violation.} The known-channel oracle assumes conditionally independent, symmetric channels. An attack against which it does \emph{worse} than the liar-removal oracle therefore lies outside that model. Three attacks do: echo (\budgetEchoSeven\ at $f=0.7$), camouflage (\budgetCamoSeven) and the sleeper (\budgetSleeperFive\ at $f=0.5$). The echo attacker copies honest reports, so its errors are not independent given the truth. The camouflage attacker conditions on the honest vote. The sleeper is non-stationary. These are exactly \race's failure modes (\cref{sec:limits}). An evaluator with labels can therefore use the sign of the budget to check whether a new attack is inside the model class that \race\ assumes. Uninformative liars have a budget of essentially zero (\budgetUninfSeven\ at $f=0.7$), as \cref{thm:silence} predicts for the minimax attack. \race\ pays a small estimation cost there (\overRemovalUninfSeven\ points relative to the liar-removal oracle), which is the price of not knowing who is silent.

\paragraph{A negative result: conditioning on disagreement (\raceD).} Camouflage liars earn trust on questions where every honest agent agrees. A natural fix is to fit the channels only on history questions where the heard answers disagree, falling back to all questions when fewer than five qualify. We call this \raceD. On the 14 attack settings other than camouflage, echo and sleeper, \raceD\ never differs significantly from \race\ (\raceDOtherNonTies\ non-ties in \raceDOtherCells\ cells). Against camouflage it helps at $f=0.7$ (\raceDCamoSeven\ points on average) but hurts at $f=0.3$ (\raceDCamoThree). Overall it wins \raceDCamoWins\ and loses \raceDCamoLosses\ of \raceDCamoCells\ camouflage cells. It also hurts against echo (\raceDEchoSeven\ at $f=0.7$), and against sleepers at $f=0.3$ (\raceDSleeperThree). Across all \raceDCells\ cells it wins \raceDWins\ and loses \raceDLosses. Restricting to disagreement removes the questions that pin down the honest peers' accuracies, trading one misspecification for another. We keep the frozen \race\ as the method. \raceD\ was specified and implemented before E9 was run, and \race\ was not modified afterwards (\cref{app:devlog}).

\begin{figure}[t]
\centering
\includegraphics[width=\linewidth]{figures/fig12_information_budget.pdf}
\caption{\textbf{E9.} What the liars are worth. Accuracy minus the liar-removal oracle, which knows who lies, removes them, and decodes the honest agents with their true accuracies. Violet stars: the known-channel oracle, which keeps the liars; its offset is the information budget. Filled blue: \race. Hollow blue: \raceD. Shaded: better than removing every liar. Means over six benchmarks and three seeds (TEST split). The sleeper is omitted because a history-fitted oracle is invalid by construction; see \cref{tab:budget}.}\label{fig:budget}
\end{figure}
